\section{Derivatives in separated cells}
\label{subsec:interacting-separated}

To estimate mixed logarithmic derivatives, we use the
real covariance $G=D_*(\mathscr A)^{-1}$. On the
coarse Grassmann algebra, write an even element
without constant term as
\[
 F=\sum_{p\ge2}\frac1{p!}\sum_{u_1,\ldots,u_p}
 F_p(u_1,\ldots,u_p)\Psi(u_1)\cdots\Psi(u_p),
\]
with antisymmetric kernels.
For a symmetric weight $w\ge1$, define
\begin{align}
 \rho_G(F;w)&=\sum_{p\ge2}\frac{h_G^p}{(p-1)!}
       \sup_{u_1}\sum_{u_2,\ldots,u_p}|F_p(u_1,\ldots,u_p)|w,
                                                     \notag\\
 N_G(F;w)&=\sum_{p\ge2}\frac{h_G^p}{(p-1)!}
       \sum_{u_1,\ldots,u_p}|F_p(u_1,\ldots,u_p)|w .
                                     \label{eq:interacting-tree-norms}
\end{align}
The supremum in $\rho_G$ is taken separately
at each degree. The indices include site, spin,
and barred or unbarred status; their distance
is the distance between the sites. We use
\[
 w_\delta(u_1,\ldots,u_p)=
          \mathrm e^{\vartheta_0\max_{i,j}d(u_i,u_j)},\qquad
 w_P(u_1,\ldots,u_p)=\mathrm e^{\vartheta_0\max_i d(u_i,P)} .
\]
The constants $h_G,\vartheta_0,\alpha_G$ are
those of \cite[(74)]{TerminalPressure}.
Antisymmetrizing a kernel that is not antisymmetric
bounds $N_G$ by the same expression for that
kernel. For $\rho_G$, replace the supremum at
the first argument by the largest supremum
over all argument positions.

In the fine algebra absorb the terminal mass
into each single-cube activity:
\[
 E^a(X;t)=\mathscr E_*(X;\mathscr A,\mathscr C(t))
              -\mathbf1_{\{X\text{ is an }M\text{-cube}\}}
                         m_*\langle\bar\psi,\psi\rangle_X.
\]
Let $E^a_{nm}(X;t)$ denote its fine spinor/pair
kernels. The antisymmetric mass kernel has
norm $8M^3$: its two fine-index weights multiply
$48\eta_*^{-3}$, its diagonal Kronecker factors
leave $48(M/\eta_*)^3$ chamber sites, and there
are four spin components and two orders.
It follows from \eqref{eq:interacting-activity-input}
that, when $a_w\le e_*^{-1/4+\varepsilon}$ and
$(4+M^3/2)e_*^\varepsilon\le1$,
\begin{equation}\label{eq:interacting-absorbed-activity}
 \sum_{n,m}\frac{a_w^{n+m}}{n!m!}\|E^a_{nm}(X;t)\|
          \le w_*^{\rm E}\mathrm e^{-\kappa d_M(X)},\qquad
 w_*^{\rm E}=e_*^{1/4-7\varepsilon}
                                      \quad(|t|_\infty\le t_x).
\end{equation}
Cauchy's formula on the circles of radius $t_x/2$
in one or two distinct mark variables gives,
on $|t|_\infty\le t_x/2$, the same bound
for $\partial_cE^a_{nm}$ and
$\partial_c\partial_{c'}E^a_{nm}$ multiplied
respectively by $q_x=2/t_x$ and $q_x^2$.
The circle integrals precede the finite weighted
kernel sums, so the coefficient norm is bounded
without replacing it by a supremum of individual
coefficients. Locality gives
$\partial_cE^a(X;t)=0$ unless $X\cap P_c\ne\varnothing$.

After substituting the terminal field maps, put
\begin{align}
 V(t)&=-\langle\bar\Psi,K(t)\Psi\rangle
           +\sum_X E^a(X;t)[H_*^\#(\mathscr A+\mathscr C(t))\Psi],
                                                        \notag\\
 K(t)&=D_*(\mathscr A+\mathscr C(t))-D_*(\mathscr A),\qquad
 V_0(t)=\sum_XE^a_{00}(X;t),\qquad V_+(t)=V(t)-V_0(t).
                                    \label{eq:interacting-tree-interaction}
\end{align}
It is even and analytic, and $V_+$ has no
constant term. Gaussian integration gives
\begin{equation}\label{eq:interacting-real-Gaussian-factor}
 Z_{\rm tm}(t)=\det D_*(\mathscr A)\mathrm e^{V_0(t)}
                                  \int\dd\mu_G\,\mathrm e^{V_+(t)} .
\end{equation}
The choice of the real covariance will allow
the exponential weights in the tree expansion.

\begin{proposition}\label{prop:interacting-differentiated-norms}
Suppose
\begin{equation}\label{eq:interacting-tree-smallness}
 5h_{\rm L}h_G\le e_*^{-1/4+\varepsilon},\qquad
 (4+M^3/2)e_*^\varepsilon\le1,\qquad
 \alpha_G\bar r\le1/2,
 \quad
 \bar r=4h_G^2C_Ke_*^{2\varepsilon}
                            +650K_0e_*^{1/4-7\varepsilon}.
\end{equation}
Let $d_{\rm sep}=\min_{c\ne c'}d(P_c,P_{c'})$ and
suppose $d_{\rm sep}\ge6M+8$.
With
$s_F=(24/\gamma^F)^3$ and
$C_V^{\mathrm{ins}}=10^3C_H^2s_F^4K_0$, we have
\begin{align}
 \rho_G(V_+(t);w_\delta)&\le\bar r
                                      &&(|t|_\infty\le t_x),\notag\\
 N_G(\partial_cV_+(t);w_{P_c})&\le q_xC_V^{\mathrm{ins}}h_G^2p_\beta,
 \quad
 \rho_G(\partial_cV_+(t);w_{P_c})\le q_xC_V^{\mathrm{ins}}h_G^2,\notag\\
 N_G(\partial_c\partial_{c'}V_+(t);1)
          +|\partial_c\partial_{c'}V_0(t)|
 &\le q_x^2C_V^{\mathrm{ins}}h_G^2p_\beta
                    \mathrm e^{-\vartheta_0(d_{\rm sep}-4M-6)} .
                             \label{eq:interacting-differentiated-norms}
\end{align}
The last two lines hold on $|t|_\infty\le t_x/2$,
with $c\ne c'$ in the last line. If the lattice
probe vanishes, all marked derivatives vanish.
\end{proposition}

\begin{proof}
We give the kernel summation which preserves
the marked supports. For a polymer $X$, let
$X^+$ be its unit-site one-neighborhood.
Every fine index in an activity supported on
$X$ has its position set in $X^+$. This
includes both chamber sites in an internal
pair and the first site of an external
nearest-neighbor pair whose other endpoint
lies over $X$. The polymer diameter estimate
gives
\begin{equation}\label{eq:interacting-X-diameter}
 \vartheta_0\operatorname{diam}(X^+)\le d_M(X)+3 .
\end{equation}
The weighted kernel after substituting the
field maps is, at degree $p=n+m$,
\[
 K_{nm}^X(u_1,\ldots,u_p)=
 \int E^a_{nm}(X;t;v_1,\ldots,v_p)
           \prod_{i=1}^p\operatorname{Leg}(v_i;u_i)\,\dd v .
\]
Here $\dd v$ includes the fine weights already
specified. Passing from the bidegree expression
to a degree-$p$ kernel produces the factor
$p!/(n!m!)$; in the tree norms its remaining
factor is $p h_G^p/(n!m!)$.

Suppose a weight is bounded by
$c_w\prod_i \mathrm e^{\vartheta_0d(u_i,\operatorname{pos}(v_i))}$.
Each summed leg contributes at most $h_{\rm L}$.
If one argument $u$ is fixed, its leg contributes
at most
$h_{\rm L}\mathrm e^{-(\gamma^F/4)d(u,X^+)}$.
Thus the full kernel sum is at most
$c_wh_{\rm L}^p\|E^a_{nm}(X;t)\|$, and the
sum with one fixed argument retains the
extra exponential in $d(u,X^+)$.
For $w_\delta$, the triangle inequality between
the two farthest arguments and their position
sets, followed by \eqref{eq:interacting-X-diameter},
permits $c_w=\mathrm e^{d_M(X)+3}$.

The needed rooted polymer sum is
\begin{equation}\label{eq:interacting-rooted-polymer-sum}
 \sum_X \mathrm e^{-(\gamma^F/4)d(u,X^+)}
                       \mathrm e^{-(\kappa-1)d_M(X)}\le32K_0.
\end{equation}
Indeed the sum of
$\mathrm e^{-(\gamma^F/4)d(u,\square^+)}$ over $M$-cubes
is at most $32$. The neighboring $27$ cubes
contribute at most $27$; a cube $j$ layers
away has distance at least $(j-1)M$ and there
are at most $(2j+1)^3$ such cubes.
Use $\gamma^FM/4\ge4$ to sum the remaining
geometric tail. Assigning a cube of $X$
and applying the anchored polymer sum $K_0$
proves \eqref{eq:interacting-rooted-polymer-sum}.

Apply \eqref{eq:interacting-absorbed-activity}
with $a_w=5h_{\rm L}h_G$. The factor $5^{-p}$
absorbs the polynomial degree factors.
The preceding full and rooted sums give
\[
 \rho_G\left(\sum_X E^a(X;t)[H_*^\#\Psi]_+;w_\delta\right)
                                       \le650K_0w_*^{\rm E}.
\]
More generally, for kernels bounded by
$\mathcal B \mathrm e^{-\kappa d_M(X)}$ at that
weight, these sums give
$650K_0\mathcal B$ in $\rho_G$ and
$\exp(3)\mathcal B\sum_X\mathrm e^{-(\kappa-1)d_M(X)}$
in $N_G$. This follows by summing
$\sum_{p\ge2}p5^{-p}=9/80$ in the rooted
case and using $p\le5^p$ in the full case.
The quadratic term contributes at most
$4h_G^2C_Ke_*^{2\varepsilon}$, since
$\vartheta_0\le\gamma^F/2$. This proves the
first line of \eqref{eq:interacting-differentiated-norms}.

For the first derivative, there are three
types of term: the coarse quadratic derivative,
the activity-kernel derivative, and a derivative
of one leg. The first has respective bounds
$8h_G^2C_K\epsilon_cp_\beta$ and
$4h_G^2C_K\epsilon_c$ in $N_G$ and $\rho_G$,
with $\epsilon_c=e_*|\widehat c_c|_\infty\le q_x$.
In its chain through $z\in P_c$, the weight
$w_{P_c}$ is absorbed by the exponential
distances from the two coarse endpoints to $z$.

For the kernel derivative, $X$ must meet $P_c$.
The same weight comparison uses that meeting
point and \eqref{eq:interacting-X-diameter}.
The Cauchy bound multiplies
$\mathcal B=w_*^{\rm E}$ by $q_x$, and
the sum of polymers meeting $P_c$ is at most
$K_0p_\beta$. Its $N_G$ bound is
$\exp(3)K_0p_\beta q_xw_*^{\rm E}$ and its
$\rho_G$ bound is $650K_0q_xw_*^{\rm E}$.

For a differentiated leg at position $j$, the
insertion estimate has a chain
$w\in\operatorname{pos}(v_j)\to z\in P_c^+\to u_j$.
The weight obeys
\[
 w_{P_c}\le \mathrm e^{d_M(X)+5}
 \mathrm e^{\vartheta_0[d(w,z)+d(z,u_j)]}
            \prod_{i\ne j}\mathrm e^{\vartheta_0d(u_i,\operatorname{pos}(v_i))}.
\]
The extra two in the exponent accounts for
$d(P_c^+,P_c)\le M+1$.
Summing the differentiated leg and the other
legs leaves, before the degree sum,
\[
 2\mathrm e^{d_M(X)+5}h_{\rm L}^{p-1}
       C_H\epsilon_cs_F
       \sum_{z\in P_c^+}\mathrm e^{-(\gamma^F/2)d(z,X^+)}
                                      \|E^a_{nm}(X;t)\|.
\]
There are $p$ choices of the differentiated
leg. The kernel norm contributes one further
factor $p$, and $p^2\le5^p$ absorbs both.
Use \eqref{eq:interacting-rooted-polymer-sum}
with $\gamma^F/2$ in place of $\gamma^F/4$.
The result is bounded by
$9500C_Hs_FK_0w_*^{\rm E}p_\beta\epsilon_c$
in $N_G$.

For $\rho_G$, if the fixed argument lies on
the differentiated leg, split its two-edge
exponent as $7\gamma^F/8=3\gamma^F/8+\gamma^F/2$.
Summing the intermediate site leaves
$s_F\mathrm e^{-(\gamma^F/4)d(u,X^+)}$.
If the fixed argument lies on another leg,
that leg supplies the same rooted exponential,
and the differentiated leg contributes
$2C_H\epsilon_cs_F^2$.
Both cases therefore have the bound
$9500C_Hs_F^2K_0w_*^{\rm E}\epsilon_c$.
Here the degree sum is
$\sum_{p\ge2}p^25^{-p}=43/160$.
Adding the three types of term, using
$w_*^{\rm E}\le1/8$, $C_K\le C_Hs_F^2$,
and $C_H\ge2$, proves the second line with
$C_V^{\mathrm{ins}}=10^3C_H^2s_F^4K_0$.

For the mixed derivative, put
\[
 d_0=d(P_c^+,P_{c'}^+),\quad
 d_1=d(X^+,P_c^+),\quad d_2=d(X^+,P_{c'}^+).
\]
Then $d_0\ge d_{\rm sep}-2M-2$.
The differentiated coarse quadratic term is
bounded in $N_G$ by
\[
 16h_G^2C_H^2s_F^3p_\beta q_x^2
                                      \mathrm e^{-(\gamma^F/2)d_{\rm sep}}.
\]
For an activity kernel with $p$ legs, there
are $(p+1)^2$ placements of the two derivatives.
They may act on two legs, the same leg, a
kernel and a leg, or both on the kernel.
The summed first differentiated legs are
bounded by $2C_H\epsilon_cs_F^2\mathrm e^{-(\gamma^F/2)d_1}$
and its $c'$ version. The mixed leg is
bounded by
\[
 2C_H^2\epsilon_c\epsilon_{c'}s_F^3\mathrm e^{-(\gamma^F/2)d_0}
       \bigl(\mathrm e^{-(\gamma^F/2)d_1}+\mathrm e^{-(\gamma^F/2)d_2}\bigr).
\]
If a derivative acts on the kernel, its
corresponding distance $d_i$ is zero by
locality. Normalize the first kernel
derivatives by $q_x$ and the mixed one by
$q_x^2$. Each of the four resulting kernel
types satisfies \eqref{eq:interacting-absorbed-activity}.
Every product-rule term is consequently at most
$4q_x^2C_H^2s_F^4h_{\rm L}^p\|E_{\rm type}\|\mathcal D(X)$,
where
\[
 \mathcal D(X)=\mathrm e^{-(\gamma^F/2)(d_1+d_2)}
       +\mathrm e^{-(\gamma^F/2)d_0}
                  \bigl(\mathrm e^{-(\gamma^F/2)d_1}+\mathrm e^{-(\gamma^F/2)d_2}\bigr).
\]

The triangle inequality through $X^+$ gives
$d_0\le d_1+d_2+M d_M(X)+2M+2$.
Use $\vartheta_0\le\gamma^F/8$ and
$\vartheta_0M\le1$ to obtain
\[
 \mathcal D(X)\mathrm e^{-\kappa d_M(X)}
 \le2\mathrm e^{-\vartheta_0(d_0-2M-2)}
       \bigl(\mathrm e^{-(\gamma^F/4)d_1}+\mathrm e^{-(\gamma^F/4)d_2}\bigr)
                                \mathrm e^{-(\kappa-1)d_M(X)}.
\]
Sum the two distance factors over the sites
of $P_c^+$ and $P_{c'}^+$, then use
\eqref{eq:interacting-rooted-polymer-sum}.
The degree factor $p(p+1)^2$ is at most
$5^p$ for $p\ge2$. At $p=0$ there is
one kernel term with no legs, and the same
bound controls $\partial_c\partial_{c'}V_0$.
The four kernel types contribute in total
at most
$256C_H^2s_F^4K_0p_\beta q_x^2
\mathrm e^{-\vartheta_0(d_{\rm sep}-4M-4)}$.
Adding the quadratic term proves the last
line of \eqref{eq:interacting-differentiated-norms}.
If the lattice probe is zero, $\mathscr C(t)=0$
and there is no mark dependence.
\end{proof}

We next turn these kernel estimates into a bound
for the logarithm. The finite-dimensional tree
formula is recalled with its proof, since the
distance between its two marked vertices is
the source of the mixed-cell decay.

For $n$ copies of the generators, let
$\Delta^{ij}=\sum_{x,y}G(x,y)
\partial_{\bar\Psi^{(j)}(y)}\partial_{\Psi^{(i)}(x)}$
and $\Delta_{\{i,j\}}=\Delta^{ij}+\Delta^{ji}$.
The derivatives are left Grassmann derivatives.
For a forest $T$ and edge parameters
$\tau\in[0,1]^T$, let $s_{ij}^T(\tau)$ be the
minimum parameter on the path from $i$ to $j$,
zero if there is no path, and $s_{ii}^T=1$.
Write $\Delta[s]=\sum_{i,j}s_{ij}\Delta^{ij}$
and let $\epsilon_0$ extract the constant term.
For even elements $F_1,\ldots,F_n$ without
constant term, define
\begin{equation}\label{eq:interacting-connected-tree}
 T_n(F_1,\ldots,F_n)=
 \sum_{\text{trees }T}\int_{[0,1]^T}
 \epsilon_0\left\{\prod_{\ell\in T}\Delta_\ell
          \mathrm e^{\Delta[s^T(\tau)]}\prod_iF_i(\Psi^{(i)})\right\}\,\dd\tau .
\end{equation}
For $n=1$ this is the Gaussian expectation.

\begin{proposition}\label{prop:interacting-tree-log}
If $V_+(t)$ satisfies
\eqref{eq:interacting-tree-smallness} and is analytic,
then
\[
 W(t)=V_0(t)+\sum_{n\ge1}\frac1{n!}
                       T_n(V_+(t),\ldots,V_+(t))
\]
is analytic on $\mathbb D(t_x)$ and
$Z_{\rm tm}(t)=\det D_*(\mathscr A)\mathrm e^{W(t)}$.
For $c\ne c'$ and nonempty unit-site sets $P,P'$,
\begin{align}
 |\partial_c\partial_{c'}(W-V_0)|
 \le{}&2N_G(\partial_c\partial_{c'}V_+;1)\notag\\
 &+2\alpha_GN_G(\partial_cV_+;w_P)
       \rho_G(\partial_{c'}V_+;w_{P'})
                                  \mathrm e^{-\vartheta_0d(P,P')}.
                                  \label{eq:interacting-tree-mixed}
\end{align}
\end{proposition}

\begin{proof}
The forest interpolation formula for a polynomial $H$ is
\[
 H(\mathbf1)=\sum_F\int
       \left(\prod_{\ell\in F}\partial_{s_\ell}H\right)
                                 (s^F(\tau))\,\dd\tau.
\]
For a partition $\pi$ of the vertices, keep fixed
entries $s_{\mathrm{in}}(i,j)$ inside its blocks and
put $X_\pi(u)_{ij}=u$ on pairs joining distinct blocks.
The chain rule and the fundamental theorem of calculus give
\begin{equation}\label{eq:interacting-forest-first-step}
 G(X_\pi(u))=G(X_\pi(0))
       +\sum_{\substack{\ell=(i,j)\\
                        i,j\text{ in distinct blocks of }\pi}}
          \int_0^u(\partial_{s_\ell}G)(X_\pi(w))\,\dd w .
\end{equation}
When $\ell$ joins two blocks, merge them and assign
the value $w$ to every pair between those two blocks.
The interpolation point is unchanged. Begin with singleton
blocks, $G=H$ and $u=1$, and iterate this identity
on each integrand. Each derivative
merges two components, so after at most $n-1$
steps the procedure ends. The successive
parameters are ordered decreasingly; when two
vertices merge, their common value is the
minimum parameter on their eventual path.
For a fixed forest, its edge orderings divide
the parameter cube into these ordered simplices.
Summing them gives the formula.

The matrix $s^F$ is a Gram matrix of unit
vectors: at $u\in[0,1]$, assign to vertex $i$
the basis vector of its component among the
edges with $\tau_\ell>u$, and integrate their
scalar products in $u$. This scalar product
is the length of the interval on which the
two vertices remain connected, namely
$s_{ij}^F$. Wick's formula and the Gram
inequality therefore bound a residual monomial
of degree $p$ under $\mathrm e^{\Delta[s^F]}$
by $g^p$, where $\|G\|_2\le g^2$.
Indeed, its determinant entries are
$\langle\omega_i\otimes\delta_x,\omega_j\otimes G\delta_y\rangle$.
These facts apply whether or not the
interpolated covariance is invertible.

Let vertex $i$ of a tree have degree $k_i$
and let the chosen homogeneous part of $F_i$
have degree $p_i$. The edge derivatives leave
the factorial $(p_i-k_i)!^{-1}$ and $p_i-k_i$
uncontracted fields at that vertex; terms with
$p_i<k_i$ vanish. Wick's Gram bound supplies
$g^{p_i-k_i}$. Sum contracted arguments from
the leaves to vertex $1$, using the weighted
row and column bound $\alpha_G$ for every
edge. If nonnegative symmetric weights $w_i$
and a distance $D$ satisfy
\[
 \mathrm e^{\vartheta_0D}\le
       \prod_iw_i(\text{arguments at }i)
       \prod_{\ell=(i,j)}
          \mathrm e^{\vartheta_0d(\text{edge arguments})},
\]
this yields $\mathrm e^{-\vartheta_0D}\alpha_G^{n-1}$
times the full kernel sum at vertex $1$
and the rooted sums at the other vertices.
The number of trees with degrees $k_1,\ldots,k_n$
is $(n-2)!/\prod_i(k_i-1)!$; here $n\ge2$,
$k_i\ge1$ and $\sum_i k_i=2n-2$, as in
\cite[Proposition~4]{PostnikovLin2019}.
Enlarge the sum over degrees and use
\[
 \sum_{k=1}^{p}
       \frac{g^{p-k}}{(k-1)!(p-k)!}
       =\frac{(1+g)^{p-1}}{(p-1)!}
       \le\frac{h_G^p}{(p-1)!}.
\]
We obtain, for $n\ge2$,
\begin{equation}\label{eq:interacting-tree-bound}
 |T_n(F_1,\ldots,F_n)|
 \le \mathrm e^{-\vartheta_0D}(n-2)!\alpha_G^{n-1}
                N_G(F_1;w_1)\prod_{i=2}^n\rho_G(F_i;w_i).
\end{equation}
For $n=1$ the bound is $N_G(F_1;1)$.
This also explains why the tree norms use
$(p-1)!$ rather than $p!$.

Apply the forest formula to the replicated
Gaussian expectation of $V_+^n$.
Different forest components contain disjoint
generator sets and their even factors commute,
so their expectations factor. Grouping by
component partitions gives
\[
 \int\dd\mu_G\,V_+^n
       =\sum_{\{B_1,\ldots,B_q\}}
                     \prod_{a=1}^qT_{|B_a|}(V_+,\ldots,V_+).
\]
The exponential formula for partitions
identifies the Gaussian integral with the
exponential of the connected tree sum.
The identity first holds as a formal series
in a scalar multiplying $V_+$. Bound
\eqref{eq:interacting-tree-bound} and
$\alpha_G\rho_G(V_+;1)\le1/2$ make the
connected series absolutely convergent up
to that scalar equal to one, so the formal
identity is an analytic identity there.
The tree terms are polynomial in the
activity coefficients, and their series
converges locally uniformly in $t$. Thus $W$
is analytic and its derivatives are obtained
term by term.

Two derivatives of the symmetric multilinear
tree expression either act on the same vertex
or on two different vertices. The first
terms, including $n=1$, sum to at most
$2N_G(\partial_c\partial_{c'}V_+;1)$.
For the second terms put the two derivatives
at vertices $1$ and $2$. The weights there
are $w_P,w_{P'}$, and those at all other
vertices are $w_\delta$.
Along the unique path from $1$ to $2$,
the distance $d(P,P')$ is at most the
sum of the two endpoint distances, all
edge distances, and the argument diameters
at intermediate vertices. Hence the weight
condition for \eqref{eq:interacting-tree-bound}
holds with $D=d(P,P')$. The tree sum leaves
$\alpha_GN_G(\partial_cV_+;w_P)
\rho_G(\partial_{c'}V_+;w_{P'})
\mathrm e^{-\vartheta_0d(P,P')}$ times
$\sum_{n\ge2}(\alpha_G\rho_G(V_+;w_\delta))^{n-2}\le2$.
This proves \eqref{eq:interacting-tree-mixed}.
\end{proof}

For a function $f$ of the marks, let
$t^T$ keep the coordinates indexed by $T$
and set all others to zero, and define
its component with variable set $S$ by
\[
 f_S(t)=\sum_{T\subset S}(-1)^{|S\setminus T|}f(t^T).
\]
Then $f=\sum_Sf_S$, $f_\varnothing=f(0)$,
and $f_S$ vanishes if any variable of $S$
vanishes. The latter statement follows by
pairing the terms with and without that
variable. In particular a variable absent
from $f$ makes every component containing
it zero. Two coordinate differences can
be represented by two integrals of the
mixed derivative. Consequently
\begin{equation}\label{eq:interacting-Mobius-bound}
 \sup_{|t|_\infty\le R}|f_S(t)|
 \le2^{|S|-2}R^2
          \sup_{|t|_\infty\le R}|\partial_c\partial_{c'}f(t)|
 \qquad(c,c'\in S,\ c\ne c').
\end{equation}

\begin{proposition}\label{prop:interacting-mixed-logarithm}
Assume the factor and mass conditions
\eqref{eq:interacting-factor-conditions}--\eqref{eq:interacting-small-complex-mass},
the tree conditions \eqref{eq:interacting-tree-smallness}, and
$d_{\rm sep}\ge6M+8$. For $|S|\ge2$, the normalized logarithm satisfies
\begin{equation}\label{eq:interacting-mixed-logarithm}
 \sup_{|t|_\infty\le t_x/2}|\mathcal G_S(t;\theta,A)|
                     \le\varpi_{\rm sep}\le\varpi_{\rm sep}^+,
\end{equation}
where
\begin{align}
 C_W^{\mathrm{mix}}&=3C_V^{\mathrm{ins}}h_G^2+2\alpha_G(C_V^{\mathrm{ins}})^2h_G^4
                             +4\exp(6)C_dM^3K_0,\notag\\
 \varpi_{\rm sep}&=2^{N_*}C_W^{\mathrm{mix}}p_\beta
                    \mathrm e^{-\vartheta_0(d_{\rm sep}-4M-6)},\notag\\
 \varpi_{\rm sep}^+&=2^{N_*}C_W^{\mathrm{mix}}p_\beta
                    \mathrm e^{-\vartheta_0(Mr_*/(8L)-4M-7)} .
                                  \label{eq:interacting-separation-errors}
\end{align}
\end{proposition}

\begin{proof}
The tree representation and uniqueness of the
normalized logarithm give
\[
 \mathcal G(t)=\sum_X\mathfrak P_*(X,\mathscr C(t))+W(t)-W(0).
\]
The component of a determinant activity is zero
unless $X$ meets every $P_c$, $c\in S$.
Such a polymer has
$d_M(X)\ge d_{\rm sep}/M-2\ge1$.
Use \eqref{eq:interacting-P-bound}, sum the
polymers anchored at one of at most $p_\beta$
cubes, and keep $\mathrm e^{-2d_M(X)}$ from
$\kappa-\kappa_0\ge2$.
Its total contribution is at most
$2^{|S|}2\exp(4)C_dM^3K_0p_\beta
\mathrm e^{-\vartheta_0d_{\rm sep}}$.

Apply \eqref{eq:interacting-tree-mixed}
with $P=P_c$, $P'=P_{c'}$ and
\eqref{eq:interacting-differentiated-norms}.
The bound for $\partial_c\partial_{c'}W$ is
\[
 q_x^2p_\beta\left[
       2C_V^{\mathrm{ins}}h_G^2\mathrm e^{-\vartheta_0(d_{\rm sep}-4M-6)}
           +2\alpha_G(C_V^{\mathrm{ins}})^2h_G^4\mathrm e^{-\vartheta_0d_{\rm sep}}\right].
\]
In \eqref{eq:interacting-Mobius-bound}, take
$R=t_x/2$ and use $q_xR=1$.
Add the determinant contribution and
$|S|\le N_*$ to obtain the first inequality.
The placement estimate
$d_{\rm sep}\ge Mr_*/(8L)-1$
proves the second.
When the lattice probe is zero,
$R_{\rm tm}=1$, $\mathcal G=0$ and the
first separation error is interpreted as zero.
\end{proof}

